% chapters/theory/07.tex

\chapter{Retrieval, Tools, and Agents: Extending LLMs Beyond Their Weights}
\label{chap:rag-tools-agents}

A language model is trained on a finite dataset, at a finite point in time, and its weights are fixed at inference. This creates an unavoidable boundary between what the model can do “from memory” and what a useful system must do in the world. The moment you ask for a recent event, a live inventory check, a calculation you want guaranteed correct, or an action that changes an external system, the model’s internal knowledge is no longer sufficient. This chapter describes three families of techniques that extend an LLM beyond its weights. Retrieval-Augmented Generation (RAG) connects the model to external text corpora. Tool calling connects the model to structured APIs and computations. Agents combine both with iterative planning so the model can pursue goals rather than answer a single question.

\section{Why LLMs need external connectivity}
A pretrained model’s knowledge is bounded by the time window of its training data. If you ask about events that occurred after that cutoff, the model can only guess based on patterns, and it will often guess incorrectly. Updating the model by continuing training is not a clean solution. Continuous updates risk regressions, and operationally they create maintenance overhead because every downstream tuned variant would need to be re-derived.

A naïve workaround is to paste new information into the prompt. This does not scale. Context windows, even when large, are finite. In addition, long prompts degrade performance because irrelevant context can distract the model, a phenomenon commonly illustrated with “needle in a haystack” tests. Finally, prompt tokens cost money and latency. These constraints motivate selective strategies that inject only the information that is likely to matter.

\section{Retrieval-Augmented Generation}
RAG is the simplest external-connectivity idea: instead of asking the model to recall everything, we retrieve relevant evidence and provide it in the prompt. The model then generates an answer grounded in that evidence. The pipeline decomposes into three steps: retrieve, augment, and generate. The quality of the system is dominated by retrieval; weak retrieval yields weak answers regardless of how strong the generator is.

\subsection{Knowledge bases, chunking, and embeddings}
RAG begins by building a knowledge base. Documents are collected from sources that are relevant to the application: policy documents, product manuals, internal wikis, papers, or any domain-specific corpus. These documents are then split into chunks—typically a few hundred tokens—because retrieval is performed at the chunk level. Each chunk is mapped to an embedding, and these embeddings are stored in an index.

Chunk size and overlap are not cosmetic hyperparameters; they determine whether the retrieved chunk is interpretable and useful. If chunks are too small, they lose context and become hard to interpret. If they are too large, embeddings become less focused and retrieval becomes less precise. Overlap mitigates boundary effects by ensuring that information needed to interpret a chunk is not cut away at the previous boundary.

\subsection{Two-stage retrieval: candidate retrieval and re-ranking}
Retrieval pipelines in large-scale systems usually follow a two-stage design borrowed from search and recommendation. Candidate retrieval filters a large corpus down to a manageable set of potentially relevant candidates. This stage is optimized for recall: it should avoid missing relevant chunks, even if it includes some noise. Re-ranking then refines the candidate set using a more expensive model, optimized for precision: it should put the truly relevant items at the top.

\subsection{Candidate retrieval with semantic similarity and bi-encoders}
The dominant candidate retrieval approach embeds the query and retrieves chunks by semantic similarity. A bi-encoder architecture encodes the query independently and encodes each chunk independently. Similarity is then computed, most often with cosine similarity, and the system retrieves the nearest neighbors in embedding space. Since corpora can contain millions of chunks, approximate nearest neighbor methods are typically used to make retrieval fast without exhaustively comparing against every chunk.

This approach enables semantic retrieval: two texts can be retrieved as relevant even if they share few or no keywords, as long as their meanings are related. Sentence-level embedding models such as Sentence-BERT are designed for exactly this purpose: they train embeddings to produce high similarity for semantically related pairs and low similarity otherwise.

\subsection{Lexical retrieval and BM25}
Semantic similarity is not always sufficient. There are many queries where exact keyword overlap is essential: identifiers, product names, legal codes, error messages, and precise terminology. Lexical retrieval methods such as BM25 score documents based on token overlap between query and document and remain valuable in production. A common practical approach is hybrid retrieval: combine semantic retrieval with lexical retrieval to improve robustness across query types.

\subsection{Bridging query/document mismatch: HyDE}
Another practical mismatch arises because user queries are often short questions while documents are long prose. Embedding a short question and a long chunk with the same encoder can yield representations that are not directly comparable. HyDE addresses this by generating a hypothetical “answer-like” document from the query using an LLM, embedding that generated document, and retrieving using the hypothetical embedding. This can improve retrieval in some settings, though it is not universally beneficial and is often treated as an optional enhancement.

\subsection{Making chunks self-contained: contextual retrieval}
Chunking inevitably separates a piece of text from its broader document context. A chunk retrieved in isolation may be technically relevant but difficult to interpret without knowing what section it came from, what definitions were introduced earlier, or what constraints apply globally.

A practical remedy is contextual retrieval. For each chunk, a short context header is generated using the full document and the chunk itself, producing a compact explanation of what the chunk is about and how it fits into the document. This context is prepended to the chunk before embedding and retrieval. The result is that retrieval returns chunks that “carry” their own local context, increasing interpretability and reducing ambiguity.

Generating these context headers can be expensive if done naïvely for every chunk. One way to reduce cost is prompt caching: because many contextualization calls share the same prefix (the full document), the system can reuse cached computation for the shared prefix and only compute the incremental work for the chunk-specific portion. Many commercial APIs expose discounted pricing for cached prefixes, but the underlying idea is general: reuse repeated prefix computation.

\subsection{Re-ranking with cross-encoders}
Candidate retrieval uses independent embeddings and a cheap similarity computation. Re-ranking often benefits from letting the model compare query and chunk jointly. A cross-encoder concatenates the query and chunk and produces a relevance score using full cross-attention. This is more expensive, but it can exploit fine-grained interactions such as negation, constraints, and detailed matching that are difficult to capture with a single dot product between independent embeddings.

\subsection{Evaluating retrieval quality}
RAG performance depends on retrieval quality. Retrieval is typically evaluated with ranking metrics under labeled relevance judgments. Precision@k and recall@k extend the familiar classification notions to the top-k retrieval setting. Mean Reciprocal Rank focuses on the position of the first relevant result, which is useful when only one correct chunk is needed. NDCG discounts relevance by rank and normalizes by the ideal ranking, rewarding systems that put relevant chunks early. Embedding benchmarks such as MTEB provide standardized evaluation suites, but domain-specific corpora often require custom relevance datasets.

\section{Tool (Function) Calling}
RAG addresses unstructured knowledge. Tool calling addresses structured interaction with external systems. Many real tasks are naturally expressed as functions: query a database, compute a route, send an email, retrieve a calendar event, calculate a price, or execute code. Tool calling allows the model to delegate such tasks to deterministic systems and then translate the result back into natural language.

Tool calling is best understood as a three-stage loop. First, the model selects the appropriate tool and produces the arguments. Second, external code executes the tool and returns a structured result. Third, the model consumes this result and produces a user-facing response. Importantly, the model does not see the tool implementation; it sees the tool schema: name, description, argument types, and output shape. The quality of schemas matters. A clear name and docstring can be the difference between reliable tool use and tool hallucination.

There are two broad ways to teach tool use. One is supervised training: provide examples where the correct action is to emit a tool call, and examples where the model must synthesize a final response from a tool output. The other is prompt-based: include clear instructions describing how to choose tools, how to format calls, and how to interpret outputs. With sufficiently capable models, prompt-based tool use can work well, especially if the prompt is iteratively improved against a held-out set of tool-use scenarios.

Tool calling spans several use cases. It can provide fresh information (search, news, prices) where a model’s static knowledge would fail. It can provide guaranteed computation (calculators, code execution) where symbolic correctness matters. It can perform actions (email, ordering, booking) where the model becomes an orchestrator rather than a source of facts.

The scaling challenge is that real systems can have many tools. If every call includes hundreds of tool definitions, the model may fail to select the right one and context budget is wasted. This motivates tool selection.

\subsection{Tool selection and routing}
Tool selection is a two-step architecture. A lightweight router receives the user query and a compact representation of the tool catalog—often just tool names and short descriptions—and selects a small subset of candidate tools. Only those tools’ full schemas are then provided to the main model, which performs tool calling. The router is typically recall-oriented: its job is not to pick the single best tool perfectly but to avoid excluding the needed tool. Tool selection can be implemented by an LLM, by a classifier, or even by retrieval-like similarity methods on tool descriptions.

\section{Standardization: Model Context Protocol}
Tool ecosystems face an interoperability problem: every LLM host and every application might define tools differently. A standard protocol makes tools portable and reduces duplication. Model Context Protocol (MCP) is one such standardization effort. The details are implementation-level, but the conceptual intent is clear: tools are served by a tool server, discovered by the host, and invoked through a consistent schema. Prompts and resources can be bundled with tools so that model hosts can use them predictably. Standardization matters because tool use becomes a platform capability rather than a one-off engineering trick.

\section{Agents: from single tool calls to goal-directed loops}
Tool calling turns an LLM into a system that can perform structured actions. Agents extend this idea by introducing iteration: instead of answering a single question, the model pursues a goal over multiple steps. The defining characteristic of an agent is a loop in which the model observes state, plans an action, executes that action via tools, observes the result, and repeats until the goal is met.

A useful conceptual framework is ReAct—reason and act—which interleaves reasoning with tool use. The vocabulary varies across implementations, but the structure is stable. The model first interprets what is currently true about the world and what is unknown. It plans the next action to reduce uncertainty or move toward the goal. It executes the action using a tool call. It then updates its internal state based on the observation and decides whether to continue.

Agents can be singular or multi-agent. In multi-agent systems, specialized agents communicate with one another, often by sending structured messages describing their capabilities and the state of ongoing tasks. Standardization efforts such as agent-to-agent protocols attempt to define how agents advertise skills, report progress, and handle cancellations.

\section{Safety and reliability in tool and agent systems}
Connecting models to tools and agents changes the risk profile. A model that can call tools can potentially access sensitive data or take irreversible actions. One concrete risk is data exfiltration: if the model can send messages or write to external channels, a malicious prompt could attempt to extract secrets. Another risk is prompt injection: an untrusted source can embed instructions that cause the model to call tools incorrectly or reveal information.

Mitigation strategies appear at two levels. Training-time mitigation includes safety data and alignment objectives that teach the model to refuse or handle unsafe requests appropriately. Inference-time mitigation adds external safeguards such as policy filters, safety classifiers, and tool-use permissions. For high-stakes tools, systems commonly require explicit user confirmation before executing actions. A separate but equally important axis is reliability. Multi-step workflows amplify error: if each step has a probability of failure, the overall probability of success declines quickly as the number of steps increases. This is one reason why fully autonomous agents remain challenging in real deployments.

\section{Practical engineering advice}
Successful systems tend to follow a disciplined development path. Start with a single narrow tool and a minimal use case and make it reliable before scaling tool catalogs. Use the most capable model available first to learn what is possible; only after correctness is established should you optimize for latency and cost. Make tool outputs small, structured, and meaningful; models often fail when tool outputs contain too much irrelevant information. Finally, treat debugging as a systematic process: failures can occur at retrieval, tool selection, argument generation, tool execution, or response synthesis, and each class of failure has different remedies.

\section{Summary}
LLMs become far more useful when they are connected to external systems. RAG extends the model’s knowledge by retrieving relevant evidence from a corpus and injecting it into the prompt. Tool calling extends the model’s capabilities by delegating structured work—retrieval, computation, or actions—to external functions and grounding the final response in deterministic outputs. Agents extend both by introducing iterative loops that plan, act, and observe, enabling goal-directed behavior. These capabilities introduce new challenges in safety and reliability, which must be addressed both through training and through inference-time safeguards.
